% SPDX-FileCopyrightText: 2025 IObundle
%
% SPDX-License-Identifier: MIT

The \texttt{VRead} and \texttt{VWrite} units (Section~\ref{sec:if_databus})
are configured with the addresses they must transfer. To avoid forcing
software to compute these addresses by hand, Versat lets the designer
describe an address-generation expression, using \texttt{for} loops and
arithmetic, that Versat turns into a generated C configuration function:

\begin{lstlisting}[language=VersatSpec,caption={An address-generation expression, and its use on \texttt{VRead}/\texttt{Mem} instances.},label={lst:addrgen_example}]
addressGen Read Simple(a){
  for x 0..a
  addr = 2 * x;
}

module AddressGenUserExample(){
  using(Simple) VRead v;
  using(Simple) Mem m;
#
  v -> m;
}
\end{lstlisting}

The \texttt{for} loop's bound (\texttt{a} above) is a runtime input to the
generated function, not a compile-time constant, so firmware can choose how
many values to transfer at run time:

\begin{lstlisting}[language=C,caption={Firmware calling the generated address-generation functions.},label={lst:addrgen_generated}]
int* input = ...;
Simple_VRead(&accelConfig->v, input, 2);
Simple_Mem(&accelConfig->m, 2);
\end{lstlisting}

As with any \texttt{for} loop in the specification language, the range
includes its start and excludes its end, so calling the generated function
with \texttt{a = 0} performs no transfer at all. This compilation step is
implemented in \texttt{software/compiler/addressGen.cpp}: each
\texttt{addressGen} definition is compiled once, at Versat compile time,
into a C function that firmware calls at run time to fill in the
associated unit's Config addresses -- so the cost of expanding the loops
into concrete addresses is paid by the CPU only when the transfer size
actually changes, not on every access.
